\itemhead{4.2 例题 · 例题2（教材p.108）}{XJ-B1-C04-EX-P108-Y0840}
\label{XJ-B1-C04-EX-P108-Y0840}
{\footnotesize 湘教版教材p.108 / PDF第115页 / 4.2 例题 / 例题2}
\field{原题}{医学中常用的钴60射线穿过1 cm厚的铅板后，强度变为原来的0.568倍。穿过厚度为 $x$ cm的铅板后的强度与原来强度之比为 $H(x)=a^x$。若铅板厚度为12 cm，射线穿过铅板后的强度与原来强度之比是多少？}
\field{参考答案}{$a=0.568,H(x)=0.568^x$；$H(12)=0.568^{12}\approx0.001128$，约为原强度的 $0.1128\%$。}
{\small \field{关联知识}{E03 指数增长、衰减与实际模型（教材帮p.131、132、154 / PDF140、141、163）}
\field{方法与易错点}{确认初始量、时间单位、倍率与时刻起点，再代入或逐次比较阈值。 易错：把年折旧10\%写成乘0.1；“超过”用成“达到”；时间起点错一位；把衰减至接近0当实际为0。}
\field{答案依据}{例题原解位置：湘教版教材p.108及题库原图；独立编写可编辑参考解答。答案册此部分主要收练习、习题，不伪造例题答案册位置。}}
